% Part: proof-theory
% Chapter: natural-deduction
% Section: translation-G2i

\documentclass[../../../include/open-logic-section]{subfiles}

\begin{document}

\olfileid[hi]{pt}{nat}{tgi}

\olsection{\Log{G2i} से \Log{N2i} में अनुवाद}

\Log{G2i} में अनुक्रमों के पूर्वपक्ष !!{formula}s के बहुसमुच्चय~$\Gamma$ हैं,
जबकि \Log{N2i} में अनुक्रमों के पूर्वपक्ष चिह्नित सूत्रों के ऐसे
समुच्चय~$\Gamma'$ हैं जिनमें प्रत्येक चिह्न केवल एक बार आता है। कहें कि
चिह्नित !!{formula}s का कोई समुच्चय~$\Gamma'$, बहुसमुच्चय~$\Gamma$ के
\emph{अनुरूप है}, यदि प्रत्येक !!{formula}~$!A$ के लिए~$\Gamma'$ में $x :
!A$ रूप के !!{element}s की संख्या,~$\Gamma$ में~$!A$ की बहुलता (अर्थात्~$!A$
की प्रतियों की संख्या) से कम या बराबर हो।

\begin{prop}\ollabel{prop:G2i-to-N2i}
यदि $\Log{G2i} + \Cut \Proves \Gamma \Sequent \Delta$, तो $\Log{N2i}
\Proves \Gamma' \Sequent !C$, जहाँ $\Delta$ रिक्त होने पर $!C \ident
\lfalse$ और अन्यथा $\Delta = \{!A\}$, तथा $\Gamma'$,~$\Gamma$ के अनुरूप है।
\end{prop}

\begin{proof}
  हम दिखाते हैं कि यदि $\pi$,~\Log{G2ci} में $\Gamma \Sequent \Delta$ का
  कोई !!a{proof} है, तो~$\Gamma$ के अनुरूप कोई~$\Gamma'$ और~\Log{N2i} में
  ~ $\Gamma' \Sequent !C$ का कोई !!a{proof}~$\delta$ है। यह हम $n =
  \pheight{\pi}$ पर आगमन से करते हैं।

  यदि $n = 0$, तो $\pi$ कोई अभिगृहीत है। यदि वह अभिगृहीत $\lfalse \Sequent
  \quad$ है, तो $\delta$ अभिगृहीत $x:\lfalse \Sequent \lfalse$ है, अर्थात्
  $\Gamma' = \{x:\lfalse\}$ और $!C \ident \lfalse$। यदि $\pi$, $!A \Sequent
  !A$ रूप का अभिगृहीत है, तो $\delta$, $x: !A \Sequent !A$ है।
  
  यदि $n>0$, तो $\pi$ किसी नियम~$R$ पर समाप्त होता है। आगमन परिकल्पना
  प्रत्याभूत करती है कि~\Log{N2c} में उस नियम के आधार-वाक्यों के अनुरूप
  अनुक्रमों के !!{proof}s हैं। हम मान सकते हैं कि उन प्रमाणों के सभी
  मुक्तिकरण चिह्न भिन्न हैं और कोई चिह्न $x$, यदि कहीं मुक्त होता है, तो केवल
  उसे मुक्त करने वाले अनुमान के ऊपर आता है (आवश्यक होने पर पुनः चिह्नित करके)।
  फिर हम दिखाते हैं कि इन !!{proof}s को मिलाकर उपयुक्त $\Gamma' \Sequent !C$
  के प्रमाण कैसे दिए जा सकते हैं। हम~$R$ के अनुसार स्थितियों में भेद करते हैं।

  \begin{enumerate}
    \item यदि $R$, \RightR{\Weakening} है, तो $\pi$ का अंत है
    \[
    \AxiomC{}
    \RightLabel{$\pi_1$}
    \Deduce$\Gamma \fCenter$
    \RightLabel{\RightR{\Weakening}}
    \UnaryInf$\Gamma \fCenter !A$
    \DisplayProof
    \]
    $\pi_1$ पर लागू आगमन परिकल्पना~$\Gamma' \Sequent \lfalse$ का कोई
    !!a{proof} देती है। $\Gamma' \Sequent !A$ का !!a{proof} पाने के लिए हम
    $\FalseInt$ लागू करते हैं।
    \item यदि $R$, \LeftR{\Weakening} है, तो $\pi$ का अंत है
    \[
    \AxiomC{}
    \RightLabel{$\pi_1$}
    \Deduce$\Gamma \fCenter \Delta$
    \RightLabel{\RightR{\Weakening}}
    \UnaryInf$!A, \Gamma \fCenter \Delta$
    \DisplayProof
    \]
    $\pi_1$ पर लागू आगमन परिकल्पना~$\Gamma' \Sequent !C$ का कोई !!a{proof}
    देती है, जिसे हम $\delta$ लेते हैं। ध्यान दें कि यदि $\Gamma'$, $\Gamma$
    के अनुरूप है, तो वह $!A, \Gamma$ के भी अनुरूप है, क्योंकि $\Gamma'$ में
    $x : !A$ की संख्या~$\Gamma$ में~$!A$ की प्रतियों की संख्या से अधिक नहीं,
    और इसलिए $\Gamma \cup \{!A\}$ में $!A$ की प्रतियों की संख्या से भी अधिक
    नहीं है।
    \item यदि $R$, \LeftR{\Contraction} है, तो $\pi$ का अंत है
    \[
    \AxiomC{}
    \RightLabel{$\pi_1$}
    \Deduce$!A, !A, \Gamma \fCenter \Delta$
    \RightLabel{\LeftR{\Contraction}}
    \UnaryInf$!A, \Gamma \fCenter \Delta$
    \DisplayProof
    \]
    $\pi_1$ पर लागू आगमन परिकल्पना~$\Pi, \Gamma' \Sequent !C$ का कोई
    !!a{proof} देती है, जहाँ $\Pi \subseteq \{x : !A, y:!A\}$। यदि $\Pi =
    \{x : !A, y:!A\}$, तो~$\pi_1$ के सभी चिह्नित !!{formula}s $y:!A$ को
    $x:!A$ से बदलें। चूँकि $x$,~$\pi_1$ के अंतिम अनुक्रम के प्रसंग में आता है,
    हम मान सकते हैं कि~$\pi_1$ का कोई अनुमान~$x:!B$ रूप का !!a{formula}
    मुक्त नहीं करता, अर्थात्~$\pi_1$ में किसी अनुक्रम के बाईं ओर कोई $x:!B$
    सक्रिय नहीं है। फलतः परिणाम अब भी~\Log{N2i} में !!a{proof} है।
    \item यदि $R$, $\RightR{\lif}$ है, तो $\pi$ का अंत है
    \[
    \AxiomC{}
    \RightLabel{$\pi_1$}
    \Deduce$!A, \Gamma \fCenter !B$
    \RightLabel{\RightR{\lif}}
    \UnaryInf$\Gamma \fCenter !A \lif !B$
    \DisplayProof
    \]
    आगमन परिकल्पना से $x: !A, \Gamma' \Sequent !B$ या $\Gamma' \Sequent
    !B$ का कोई !!a{proof}~$\delta_1$ है, जहाँ $\Gamma'$,~$\Gamma$ के अनुरूप
    है। हम $\delta_1$ और \Intro{\lif} के एक अनुप्रयोग को $\delta$ ले सकते
    हैं।
    \item यदि $R$, \LeftR{\lif} है, तो $\pi$ का अंत है
    \[
    \AxiomC{}
    \RightLabel{$\pi_1$}
    \Deduce$\Gamma_1 \fCenter !A$
    \AxiomC{}
    \RightLabel{$\pi_2$}
    \Deduce$!B, \Gamma_2 \fCenter \Delta$
    \RightLabel{\LeftR{\lif}}
    \BinaryInf$!A \lif !B, \Gamma_1, \Gamma_2 \fCenter \Delta$
    \DisplayProof
    \]
    आगमन परिकल्पना $\Gamma_1' \Sequent !A$ का !!{proof} $\delta_1$ और $x:
    !B, \Gamma_2' \Sequent !C$ अथवा $\Gamma_2' \Sequent !C$ का !!{proof}
    $\delta_2$ देती है। यदि $\delta_2$ दूसरे रूप का है, तो $\delta_2$ स्वयं
    ~ $\delta$ हो सकता है। अन्यथा हम~$\pi_1$ के सभी सूत्रों और मुक्तिकरण
    चिह्नों को पुनः चिह्नित करके सुनिश्चित कर सकते हैं कि कोई चिह्न
    $\delta_1$ और~$\delta_2$ दोनों में न आए। तब $\Gamma_1' \cup \Gamma_2'$,
    $\Gamma_1 \cup \Gamma_2$ के अनुरूप है। चूँकि $x:!B$,~$\delta_2$ के अंतिम
    अनुक्रम में आता है, इसलिए $\delta_2$ में मुक्तिकरण चिह्न~$x$ वाले किसी
    अनुमान में $x:!B$ सक्रिय नहीं हो सकता।~$\delta_2$ के प्रत्येक अभिगृहीत
    $x:!B \Sequent !B$ को !!{proof}~$\delta_3$ से बदलें
    \[
    \Axiom$x: !A \lif !B \fCenter !A \lif !B$
    \AxiomC{}
    \RightLabel{$\delta_1$}
    \Deduce$\Gamma_1' \fCenter !A$
    \RightLabel{\Elim{\lif}}
    \BinaryInf$x: !A \lif !B \fCenter !B$
    \DisplayProof
    \]
    और~$\delta_2$ में $x:!B$ की प्रत्येक उपस्थिति को $x:!A \lif !B$ से
    बदलें। परिणाम $\Subst{\delta_2}{\delta_3}{x:!B}$,
    $x:!A \lif !B, \Gamma_1', \Gamma_2' \Sequent !C$ का !!a{proof} है।
    \item यदि $R$, \RightR{\land} है, तो $\pi$ का अंत है
    \[
    \AxiomC{}
    \RightLabel{$\pi_1$}
    \Deduce$\Gamma_1 \fCenter !A$
    \AxiomC{}
    \RightLabel{$\pi_2$}
    \Deduce$\Gamma_2 \fCenter !B$
    \RightLabel{\RightR{\land}}
    \BinaryInf$\Gamma_1, \Gamma_2 \fCenter !A \land !B$
    \DisplayProof
    \]
    आगमन परिकल्पना $\Gamma_1' \Sequent !A$ का प्रमाण $\delta_1$ और
    $\Gamma_2' \Sequent !A$ का $\delta_2$ देती है।~$\pi_2$ में सभी सूत्रों
    और मुक्तिकरण चिह्नों को पुनः चिह्नित करके हम सुनिश्चित कर सकते हैं कि कोई
    चिह्न $\delta_1$ तथा $\delta_2$ दोनों में न आए। तब $\Gamma_1' \cup
    \Gamma_2'$, $\Gamma_1 \cup \Gamma_2$ के अनुरूप है। $\delta_1$
    और~$\delta_2$ के अंतिम अनुक्रमों पर $\Intro{\land}$ लागू करके हमें
    !!a{proof}~$\delta$ मिलता है।
    \item यदि $R$, $\LeftR{\land}$ है, तो उदाहरणतः $\pi$ का अंत है
    \[
    \AxiomC{}
    \RightLabel{$\pi_1$}
    \Deduce$!A, \Gamma \fCenter \Delta$
    \RightLabel{\LeftR{\land}}
    \UnaryInf$!A \land !B, \Gamma \fCenter \Delta$
    \DisplayProof
    \]
    आगमन परिकल्पना से $x: !A, \Gamma' \Sequent !C$ या $\Gamma' \Sequent
    !C$ का कोई !!a{proof}~$\delta_1$ है, जहाँ $\Gamma'$,~$\Gamma$ के अनुरूप
    है। दूसरी स्थिति में हम $\delta_1$ को $\delta$ ले सकते हैं। पहली स्थिति
    में ध्यान दें कि $x:!A$ अंतिम अनुक्रम में आता है, इसलिए $\delta_1$ में
    मुक्तिकरण चिह्न~$x$ वाले किसी अनुमान में $x:!A$ सक्रिय नहीं है। प्रत्येक
    अभिगृहीत $x:!A \Sequent !A$ को !!{proof}~$\delta_2$ से बदलें
    \[
    \Axiom$x: !A \land !B \fCenter !A \land !B$
    \RightLabel{\Elim{\land}}
    \UnaryInf$x: !A \land !B \fCenter !A$
    \DisplayProof
    \]
    और $x:!A$ की प्रत्येक उपस्थिति को $x:!A \land !B$ से बदलें, अर्थात्
    $\Subst{\delta_1}{\delta_2}{x:!A}$ पर विचार करें। यह $x:!A \land !B,
    \Gamma' \Sequent !C$ का !!a{proof} है।
    \item यदि $R$, \Cut है, तो $\pi$ का अंत है
    \[
    \AxiomC{}
    \RightLabel{$\pi_1$}
    \Deduce$\Gamma_1 \fCenter !A$
    \AxiomC{}
    \RightLabel{$\pi_2$}
    \Deduce$!A, \Gamma_2 \fCenter \Delta$
    \RightLabel{\Cut}
    \BinaryInf$\Gamma_1, \Gamma_2 \fCenter \Delta$
    \DisplayProof
    \]
    आगमन परिकल्पना $\Gamma_1' \Sequent !A$ का प्रमाण~$\delta_1$ और
    $x:!A, \Gamma_2' \Sequent !C$ अथवा $\Gamma_2' \Sequent !C$ का
    प्रमाण~$\delta_2$ देती है। दूसरी स्थिति में $\delta$ को~$\delta_2$
    लें। अन्यथा $\delta_2$ के प्रत्येक अभिगृहीत $x:!A \Sequent !A$ को
    $\delta_1$ से बदलें; इससे $\Gamma_1', \Gamma_2' \Sequent !C$ का
    !!a{proof}~$\delta$, अर्थात्~$\Subst{\delta_2}{\delta_1}{x:!A}$,
    प्राप्त होता है।
  \end{enumerate}
\end{proof}

\end{document}
